% =================================================
% Методическая версия для учителя — скорректированная версия
% =================================================
\worksheettitle
  {Логарифмическое неравенство}
  {Методическая версия: нотация, модуль, две замены и обратный переход}

\begin{problembox}
  Решите неравенство
  \[
    \lg^4((x^2-4)^2)-\lg^2((x^2-4)^4)\ge192.
  \]
\end{problembox}

\begin{teacherbox}[Входная диагностика: 4 минуты]
  Ответы к четырём заданиям в начале ученического листа:
  \[
    \lg^2 100=(\lg100)^2=4,
    \qquad x\ne1,
  \]
  \[
    y\in(-\infty,-3]\cup[3,+\infty),
    \qquad
    x\in(-\infty,-4]\cup[4,+\infty).
  \]
  Одна ошибка требует короткого адресного напоминания. При двух и более
  ошибках соответствующие переходы лучше выполнить на доске вместе с классом.
  Диагностика не оценивается: она определяет необходимый объём опоры.
\end{teacherbox}

\begin{notebox}[Однозначная расшифровка условия]
  В условии
  \[
    \lg^4 A=(\lg A)^4,
    \qquad
    \lg^2 A=(\lg A)^2.
  \]
  Следовательно, исходная запись означает
  \[
    \bigl[\lg((x^2-4)^2)\bigr]^4
    -\bigl[\lg((x^2-4)^4)\bigr]^2\ge192.
  \]
\end{notebox}

\section{Дидактическая карта}

\begin{center}
  \renewcommand{\arraystretch}{1.5}
  \begin{tabularx}{\textwidth}{@{}
    >{\bfseries\sffamily}p{0.29\textwidth}X@{}}
    \toprule
    \rowcolor{TableHeader}
    Параметр & Рекомендация \\
    \midrule
    Уровень & 10--11 класс, углублённая группа или математический кружок \\
    Формат & Два урока по 45 минут или одно занятие 80--90 минут \\
    Центральная тема & Различение \(\lg^2 A\), \(\lg(A^2)\), \(\log_2 A\); корректное преобразование чётных степеней с модулем \\
    Основной метод & \(y=\lg|x^2-4|\), затем \(t=y^2=\lg^2|x^2-4|\) \\
    Основной результат & Ученик сворачивает сложное выражение до квадратного неравенства и равносильно разворачивает решение обратно \\
    Основные риски & Потеря модуля, неверное чтение нотации, потеря ветви \(y\le-2\), включение \(x=\pm2\) \\
    \bottomrule
  \end{tabularx}
\end{center}

\begin{teacherbox}[Методический замысел]
  В этой версии комплект не избегает записи \(\lg^2 A\), а учит правильно
  её читать. Сначала краткая запись расшифровывается, затем обе формы
  используются параллельно:
  \[
    t=\lg^2|x^2-4|=\bigl(\lg|x^2-4|\bigr)^2.
  \]
  После закрепления допустимо пользоваться только стандартной краткой формой.
\end{teacherbox}

\notationtrianglefigure
\logchainteacherfigure

\section{Рекомендуемый сценарий занятия}

\begin{teacherbox}[Этап 1. Чтение нотации — 6 минут]
  Используйте результат входной диагностики и кратко сопоставьте три записи:
  \[
    \lg^2 A,
    \qquad
    \lg(A^2),
    \qquad
    \log_2 A.
  \]
  Если первое диагностическое задание выполнено верно, не задерживайтесь на
  дополнительных числовых примерах. Если возникла ошибка, попросите ученика
  словами указать, что является основанием, аргументом и степенью значения.

  Ожидаемое заполнение таблицы:
  \begin{itemize}[itemsep=0.2em,topsep=0.25em]
    \item \(\lg^2 A\) --- «квадрат десятичного логарифма \(A\)»;
      число \(2\) является показателем степени значения логарифма;
    \item \(\lg(A^2)\) --- «десятичный логарифм числа \(A^2\)»;
      число \(2\) является показателем степени аргумента;
    \item \(\log_2 A\) --- «логарифм \(A\) по основанию \(2\)»;
      число \(2\) является основанием логарифма.
  \end{itemize}
\end{teacherbox}

\begin{teacherbox}[Этап 2. ОДЗ и появление модуля — 15 минут]
  Проведите рассуждение в четыре шага: пример \(x=0\), сравнение значений
  при \(z=10\) и \(z=-10\), вывод общей формулы и её применение к задаче.
  Вопросы классу:
  \begin{enumerate}
    \item Почему квадрат не гарантирует положительность аргумента логарифма?
    \item Почему \(\lg((x^2-4)^2)=2\lg(x^2-4)\) неверно на всей ОДЗ?
    \item Как исправить формулу, не сужая ОДЗ?
  \end{enumerate}
  Ожидаемый результат:
  \[
    x\ne\pm2,
    \qquad
    \lg((x^2-4)^2)
      =\lg(|x^2-4|^2)
      =2\lg|x^2-4|.
  \]
\end{teacherbox}

\begin{teacherbox}[Этап 3. Первая замена — 10 минут]
  Не сообщайте сразу итоговое биквадратное неравенство. После замены
  \[
    y=\lg|x^2-4|
  \]
  предложите заполнить таблицу:
  \[
    \lg((x^2-4)^2)=2y,
    \qquad
    \lg((x^2-4)^4)=4y.
  \]
  Ученики должны самостоятельно получить
  \[
    (2y)^4-(4y)^2\ge192
    \Longleftrightarrow
    y^4-y^2-12\ge0.
  \]
\end{teacherbox}

\begin{teacherbox}[Этап 4. Вторая замена и возврат — 20 минут]
  Положите
  \[
    t=y^2=\lg^2|x^2-4|,
    \qquad t\ge0.
  \]
  После решения \((t-4)(t+3)\ge0\) оставьте на доске цепочку:
  \[
    t\ge4
    \Rightarrow y^2\ge4
    \Rightarrow |y|\ge2
    \Rightarrow y\le-2\ \text{или}\ y\ge2.
  \]
  Формулировка «извлекаем квадратный корень» нежелательна: она провоцирует
  потерю отрицательной ветви.
\end{teacherbox}

\begin{teacherbox}[Этап 5. Два режима аргумента — 25 минут]
  После возврата к логарифму класс получает
  \[
    0<|x^2-4|\le10^{-2}
    \quad\text{или}\quad
    |x^2-4|\ge10^2.
  \]
  Полезно распределить ветви между парами. После решения обсудить, почему
  ответ состоит из четырёх коротких промежутков около \(\pm2\) и двух
  далёких лучей.
\end{teacherbox}

\section{Полное математическое решение}

\stepheading{1}{Область допустимых значений}

Аргументы обоих логарифмов положительны тогда и только тогда, когда
\[
  x^2-4\ne0.
\]
Следовательно,
\[
  x\ne-2,
  \qquad
  x\ne2.
\]

\stepheading{2}{Преобразуем логарифмы без сужения ОДЗ}

На ОДЗ \(|x^2-4|>0\), поэтому
\[
  \lg((x^2-4)^2)=2\lg|x^2-4|,
\]
\[
  \lg((x^2-4)^4)=4\lg|x^2-4|.
\]
Вводим
\[
  y=\lg|x^2-4|.
\]
Тогда исходное неравенство принимает вид
\[
  (2y)^4-(4y)^2\ge192,
\]
то есть
\[
  16y^4-16y^2\ge192,
\]
\[
  y^4-y^2-12\ge0.
\]

\stepheading{3}{Решаем алгебраическое ядро}

Положим
\[
  t=y^2=\lg^2|x^2-4|,
  \qquad t\ge0.
\]
Получаем
\[
  t^2-t-12\ge0,
\]
\[
  (t-4)(t+3)\ge0.
\]
Отсюда
\[
  t\le-3
  \quad\text{или}\quad
  t\ge4.
\]
С учётом \(t\ge0\):
\[
  t\ge4.
\]
Следовательно,
\[
  y^2\ge4
  \quad\Longleftrightarrow\quad
  |y|\ge2
  \quad\Longleftrightarrow\quad
  y\le-2\ \text{или}\ y\ge2.
\]

\stepheading{4}{Возвращаемся к аргументу логарифма}

Так как \(y=\lg|x^2-4|\), получаем:
\[
  \lg|x^2-4|\le-2
  \quad\text{или}\quad
  \lg|x^2-4|\ge2.
\]
Десятичный логарифм возрастает, поэтому
\[
  0<|x^2-4|\le10^{-2}=\frac1{100}
  \quad\text{или}\quad
  |x^2-4|\ge10^2=100.
\]

\argregimesteacherfigure

\stepheading{5}{Решаем ветвь малого аргумента}

\[
  0<|x^2-4|\le\frac1{100}
\]
равносильно совокупности
\[
  -\frac1{100}\le x^2-4<0
  \quad\text{или}\quad
  0<x^2-4\le\frac1{100}.
\]
Отсюда
\[
  \frac{399}{100}\le x^2<4
  \quad\text{или}\quad
  4<x^2\le\frac{401}{100}.
\]
Первая часть даёт
\[
  \frac{\sqrt{399}}{10}\le|x|<2,
\]
то есть
\[
  x\in
  \left(-2,-\frac{\sqrt{399}}{10}\right]
  \cup
  \left[\frac{\sqrt{399}}{10},2\right).
\]
Вторая часть даёт
\[
  2<|x|\le\frac{\sqrt{401}}{10},
\]
то есть
\[
  x\in
  \left[-\frac{\sqrt{401}}{10},-2\right)
  \cup
  \left(2,\frac{\sqrt{401}}{10}\right].
\]

\stepheading{6}{Решаем ветвь большого аргумента}

\[
  |x^2-4|\ge100
\]
равносильно
\[
  x^2-4\le-100
  \quad\text{или}\quad
  x^2-4\ge100.
\]
Первое неравенство не имеет решений, поскольку \(x^2-4\ge-4\). Поэтому
\[
  x^2\ge104,
\]
\[
  |x|\ge2\sqrt{26},
\]
\[
  x\in(-\infty,-2\sqrt{26}]\cup[2\sqrt{26},+\infty).
\]

\stepheading{7}{Объединяем промежутки}

\[
\boxed{
\begin{aligned}
  x\in{}&(-\infty,-2\sqrt{26}]\\
  &\cup\left[-\frac{\sqrt{401}}{10},-2\right)
  \cup\left(-2,-\frac{\sqrt{399}}{10}\right]\\
  &\cup\left[\frac{\sqrt{399}}{10},2\right)
  \cup\left(2,\frac{\sqrt{401}}{10}\right]\\
  &\cup[2\sqrt{26},+\infty).
\end{aligned}}
\]

\finalnumberlineteacherfigure

\section{Альтернативное решение}

\begin{methodintro}[Замена всего логарифма квадрата]
  Можно положить
  \[
    s=\lg((x^2-4)^2).
  \]
  Тогда
  \[
    \lg((x^2-4)^4)=2s
  \]
  и возникает неравенство
  \[
    s^4-4s^2-192\ge0.
  \]
  Этот путь корректен, однако он не отрабатывает запись \(\lg^2 A\) и
  требует менее естественных чисел \(-12\) и \(16\) после второй замены.
  Поэтому в ученическом листе он оставлен как дополнительный, а не основной.
\end{methodintro}

\section{Типичные ошибки и диагностические вопросы}

\begin{center}
  \renewcommand{\arraystretch}{1.45}
  \begin{tabularx}{\textwidth}{@{}
    >{\raggedright\arraybackslash}p{0.35\textwidth}
    >{\raggedright\arraybackslash}X@{}}
    \toprule
    \rowcolor{TableHeader}
    Ошибка & Вопрос учителя \\
    \midrule
    \(\lg^2 A=\log_2 A\) & Где расположен показатель степени, а где основание логарифма? \\
    \(\lg^2 A=\lg(A^2)\) & Что возводится в квадрат: значение функции или её аргумент? \\
    \(\lg(z^2)=2\lg z\) при \(z<0\) & На какой области определена правая часть? Где должен появиться модуль? \\
    Из \(t=y^2\) следует \(y\ge0\) & Неотрицательна новая переменная \(t\) или исходное выражение \(y\)? \\
    Из \(y^2\ge4\) следует только \(y\ge2\) & Какие числа на числовой прямой удалены от нуля не меньше чем на 2? \\
    \(|x^2-4|\le1/100\) включает \(x=\pm2\) & Какое строгое условие было до раскрытия модуля? \\
    \bottomrule
  \end{tabularx}
\end{center}

\begin{teacherbox}[Ключевой вопрос для рефлексии]
  Почему переменная \(t\) обязана быть неотрицательной, а переменная \(y\)
  может быть отрицательной?

  \textbf{Ожидаемый ответ:} \(t=y^2\ge0\), а
  \(y=\lg|x^2-4|\) отрицателен при \(0<|x^2-4|<1\).
\end{teacherbox}

\section{Дифференциация}

\begin{teacherbox}[Минимальный маршрут]
  Оставьте готовыми формулы с модулем и первую замену. Ученики самостоятельно
  выполняют вторую замену, обратный переход и одну из двух ветвей для \(x\).
\end{teacherbox}

\begin{teacherbox}[Основной маршрут]
  Ученик заполняет блок нотации, выводит формулы с модулем, выполняет обе
  замены и решает обе ветви по таблицам рабочего листа.
\end{teacherbox}

\begin{teacherbox}[Углубление]
  Предложите сравнить два способа замены и объяснить, почему основной способ
  приводит к более простому ядру \(t^2-t-12\ge0\). Затем обобщите задачу на
  основание \(b>1\) и выражение \(\log_b|f(x)|\).
\end{teacherbox}

\newpage
\section{Критерии оценивания полного решения}

\begin{center}
  \renewcommand{\arraystretch}{1.45}
  \begin{tabularx}{\textwidth}{@{}
    >{\centering\arraybackslash}p{0.13\textwidth}
    >{\raggedright\arraybackslash}X
    >{\centering\arraybackslash}p{0.15\textwidth}@{}}
    \toprule
    \rowcolor{TableHeader}
    Этап & Что должно быть в решении & Баллы \\
    \midrule
    1 & Верно прочитана запись и найдена ОДЗ \(x\ne\pm2\) & 2 \\
    2 & Получены формулы с \(|x^2-4|\) без сужения ОДЗ & 2 \\
    3 & Введено \(y=\lg|x^2-4|\), получено \(y^4-y^2-12\ge0\) & 2 \\
    4 & Введено \(t=y^2\ge0\), получено \(t\ge4\) & 2 \\
    5 & Сохранены обе ветви \(y\le-2\) и \(y\ge2\) & 1 \\
    6 & Верно решена ветвь малого аргумента & 2 \\
    7 & Верно решена ветвь большого аргумента & 1 \\
    8 & Точно объединены промежутки и указаны границы & 2 \\
    \bottomrule
  \end{tabularx}
\end{center}

\begin{teacherbox}[Интерпретация результата]
  Максимум --- 14 баллов.
  \begin{itemize}[itemsep=0.2em,topsep=0.25em]
    \item \textbf{12--14 баллов:} метод освоен; ученик выполняет решение
      самостоятельно и обосновывает равносильность переходов.
    \item \textbf{9--11 баллов:} метод в целом освоен, но остаются отдельные
      вычислительные или оформительские неточности.
    \item \textbf{6--8 баллов:} требуется повторить обратный переход от
      переменной \(t\) к логарифму и затем к \(x\).
    \item \textbf{0--5 баллов:} требуется повторить ОДЗ, преобразование
      чётной степени с модулем и чтение логарифмической нотации.
  \end{itemize}
  Потеря модуля, потеря ветви \(y\le-2\) или включение \(x=\pm2\) считаются
  критическими содержательными ошибками, даже если итоговая сумма высока.
\end{teacherbox}
